\slajdid{10-s014}
\begin{frame}[label=10-s014]
\gusttekst\setlength{\abovedisplayskip}{3pt}\setlength{\belowdisplayskip}{3pt}\setlength{\abovedisplayshortskip}{2pt}\setlength{\belowdisplayshortskip}{2pt}I dalje važi (uvek važi)
\[\begin{gathered}V_I-R_B I_{RB}-V_{BE1}=0\\V_O=V_{CC}-R_C I_{RC}\end{gathered}\]
Za neopterećeno kolo i tranzistor u aktivnom režimi
\[V_O=V_{CC}-R_C I_{RC}=V_{CC}-R_C I_R=V_{CC}-R_C\beta_F I_B\]
Kako je
\[I_B=I_{RB}=\frac{V_I-V_{BE1}}{R_B}\]
onda je
\[V_O=V_{CC}-R_C\beta_F I_B=V_{CC}-R_C\beta_F\frac{V_I-V_{BE1}}{R_B}=V_{CC}-\beta_F\frac{R_C}{R_B}(V_I-V_{BE1})\]
Tranzistor je u aktivnom režimu pa je $V_{BE1}=V_{BE}$
\[V_O=V_{CC}-\beta_F\frac{R_C}{R_B}(V_I-V_{BE})\]
\[V_O=V_{CC}-\beta_F\frac{R_C}{R_B}(V_I-0.6V)\]
\end{frame}
